\def\tituloconcepto{Algoritmo de Bernstein--Vazirani}
% Si el capitulo se compila solo (sin main.tex), esta cabecera abre el documento.
% Antes de % Si el capitulo se compila solo (sin main.tex), esta cabecera abre el documento.
% Antes de % Si el capitulo se compila solo (sin main.tex), esta cabecera abre el documento.
% Antes de \input{cabecera} cada capitulo define \tituloconcepto.
\ifdefined\enmain\else
  \documentclass[11pt,a4paper]{article}
  \input{preambulo}
  \begin{document}
  \title{\tituloconcepto}
  \author{\autor}
  \date{\today}
  \maketitle
\fi
 cada capitulo define \tituloconcepto.
\ifdefined\enmain\else
  \documentclass[11pt,a4paper]{article}
  % Preambulo compartido por todos los documentos. Compilar con XeLaTeX/LuaLaTeX
% (o Tectonic). Si se usa Overleaf, elegir el compilador XeLaTeX.
\usepackage[spanish,es-noshorthands]{babel}
\usepackage{amsmath,amssymb}
\usepackage{booktabs,array}
\usepackage[a4paper,margin=2.5cm]{geometry}
\usepackage{xcolor}
\usepackage{listings}
\usepackage[most]{tcolorbox}
\usepackage{quantikz}
\usepackage[hidelinks]{hyperref}

% ---- Autor (editar a gusto) -------------------------------------------------
\newcommand{\autor}{Genesis Silesky}

% ---- Notacion de Dirac y atajos ---------------------------------------------
% quantikz ya define \ket y \bra; DeclareRobustCommand los redefine sin error.
\DeclareRobustCommand{\ket}[1]{\left|#1\right\rangle}
\DeclareRobustCommand{\bra}[1]{\left\langle#1\right|}
\DeclareRobustCommand{\braket}[2]{\left\langle#1\middle|#2\right\rangle}
\newcommand{\C}{\mathbb{C}}
\newcommand{\R}{\mathbb{R}}
\DeclareMathOperator{\Tr}{Tr}

% ---- Cajas de apoyo --------------------------------------------------------
% "concepto": definiciones previas que hacen falta para entender el capitulo.
\newtcolorbox{concepto}[1]{breakable,colback=blue!4,colframe=blue!55!black,
  fonttitle=\bfseries,title={Concepto previo: #1}}
% "nota": aclaraciones y advertencias.
\newtcolorbox{nota}[1]{breakable,colback=orange!6,colframe=orange!70!black,
  fonttitle=\bfseries,title={#1}}

% ---- Estilos de listados ----------------------------------------------------
\lstdefinestyle{py}{language=Python,basicstyle=\ttfamily\footnotesize,
  breaklines=true,postbreak=\mbox{\textcolor{gray}{$\hookrightarrow$}\space},
  numbers=left,numberstyle=\tiny\color{gray},frame=single,
  commentstyle=\color{green!40!black}\itshape,
  keywordstyle=\color{blue!70!black}\bfseries,
  stringstyle=\color{red!60!black},showstringspaces=false,tabsize=4,
  columns=fullflexible,keepspaces=true}
\lstdefinestyle{salida}{basicstyle=\ttfamily\footnotesize,breaklines=true,
  frame=leftline,backgroundcolor=\color{gray!8},columns=fullflexible,
  keepspaces=true}

  \begin{document}
  \title{\tituloconcepto}
  \author{\autor}
  \date{\today}
  \maketitle
\fi
 cada capitulo define \tituloconcepto.
\ifdefined\enmain\else
  \documentclass[11pt,a4paper]{article}
  % Preambulo compartido por todos los documentos. Compilar con XeLaTeX/LuaLaTeX
% (o Tectonic). Si se usa Overleaf, elegir el compilador XeLaTeX.
\usepackage[spanish,es-noshorthands]{babel}
\usepackage{amsmath,amssymb}
\usepackage{booktabs,array}
\usepackage[a4paper,margin=2.5cm]{geometry}
\usepackage{xcolor}
\usepackage{listings}
\usepackage[most]{tcolorbox}
\usepackage{quantikz}
\usepackage[hidelinks]{hyperref}

% ---- Autor (editar a gusto) -------------------------------------------------
\newcommand{\autor}{Genesis Silesky}

% ---- Notacion de Dirac y atajos ---------------------------------------------
% quantikz ya define \ket y \bra; DeclareRobustCommand los redefine sin error.
\DeclareRobustCommand{\ket}[1]{\left|#1\right\rangle}
\DeclareRobustCommand{\bra}[1]{\left\langle#1\right|}
\DeclareRobustCommand{\braket}[2]{\left\langle#1\middle|#2\right\rangle}
\newcommand{\C}{\mathbb{C}}
\newcommand{\R}{\mathbb{R}}
\DeclareMathOperator{\Tr}{Tr}

% ---- Cajas de apoyo --------------------------------------------------------
% "concepto": definiciones previas que hacen falta para entender el capitulo.
\newtcolorbox{concepto}[1]{breakable,colback=blue!4,colframe=blue!55!black,
  fonttitle=\bfseries,title={Concepto previo: #1}}
% "nota": aclaraciones y advertencias.
\newtcolorbox{nota}[1]{breakable,colback=orange!6,colframe=orange!70!black,
  fonttitle=\bfseries,title={#1}}

% ---- Estilos de listados ----------------------------------------------------
\lstdefinestyle{py}{language=Python,basicstyle=\ttfamily\footnotesize,
  breaklines=true,postbreak=\mbox{\textcolor{gray}{$\hookrightarrow$}\space},
  numbers=left,numberstyle=\tiny\color{gray},frame=single,
  commentstyle=\color{green!40!black}\itshape,
  keywordstyle=\color{blue!70!black}\bfseries,
  stringstyle=\color{red!60!black},showstringspaces=false,tabsize=4,
  columns=fullflexible,keepspaces=true}
\lstdefinestyle{salida}{basicstyle=\ttfamily\footnotesize,breaklines=true,
  frame=leftline,backgroundcolor=\color{gray!8},columns=fullflexible,
  keepspaces=true}

  \begin{document}
  \title{\tituloconcepto}
  \author{\autor}
  \date{\today}
  \maketitle
\fi


\section{Algoritmo de Bernstein--Vazirani}

\subsection{Objetivo}
Descubrir un número binario oculto de $n$ bits con una sola consulta, en vez de
las $n$ consultas que necesita un algoritmo clásico, y ver que la ventaja viene de
la superposición y no del entrelazamiento.

\subsection{Conceptos previos}
\begin{concepto}{puertas $H$ sobre varios qubits y producto punto módulo 2}
Ver capítulos~05a y~15. El \emph{producto punto} de dos cadenas de bits es
$a\cdot x=a_1x_1+a_2x_2+\cdots+a_nx_n$, y ``módulo~2'' significa quedarse con
el resto de dividir entre 2 (la paridad).
\end{concepto}

\subsection{Fundamento teórico}
El algoritmo fue propuesto por Bernstein y Vazirani en 1993. Muestra la
ventaja cuántica en términos de número de consultas.

\paragraph{El problema.} Una base de datos guarda un número oculto
$a\in\{0,1\}^n$. Se puede hacer una consulta eligiendo un $x\in\{0,1\}^n$, y la
respuesta es
\begin{equation}
  f(x)=a\cdot x\ (\mathrm{mod}\ 2)=a_1x_1\oplus a_2x_2\oplus\cdots\oplus a_nx_n,
\end{equation}
un solo bit (0 o 1). El objetivo es hallar $a$.

\paragraph{Solución clásica.} Se consulta con $x=100\ldots0$, $010\ldots0$,
$\ldots$, $000\ldots1$: la consulta $i$ da el bit $a_i$. Hacen falta \textbf{$n$
consultas}. Por ejemplo, para $a=1011$ se necesitan cuatro.

\paragraph{Solución cuántica.} Se prepara una superposición uniforme de $n$
qubits y se usa como entrada de la función, de modo que todas las consultas se
hacen en paralelo. El circuito de este ejemplo es una variante de
Phys.\ Rev.~A~\textbf{64}, 042306, y no usa qubits auxiliares:
\begin{center}
\begin{quantikz}
\lstick{$\ket0^{\otimes n}$} & \gate{H^{\otimes n}} & \gate{U_f} & \gate{H^{\otimes n}} & \meter{}
\end{quantikz}
\end{center}
con el oráculo $U_f\ket x=(-1)^{a\cdot x}\ket x$.

\paragraph{Por qué funciona.} Una $H$ sobre un qubit da
$H\ket a=\frac1{\sqrt2}\sum_{x\in\{0,1\}}(-1)^{a x}\ket x$. Para $n$ qubits:
\begin{equation}
  H^{\otimes n}\ket a=\frac1{\sqrt{2^n}}\sum_{x\in\{0,1\}^n}(-1)^{a\cdot x}\ket x ,
  \qquad
  H^{\otimes n}\ket{0}^{\otimes n}=\frac1{\sqrt{2^n}}\sum_{x}\ket x .
\end{equation}
Es decir, $U_f$ convierte la superposición uniforme en $H^{\otimes n}\ket a$.
Como $H^{\otimes n}H^{\otimes n}=I$, la $H$ final lo lleva de vuelta a
$\ket a$:
\begin{equation}
  \ket a=H^{\otimes n}\,U_f\,H^{\otimes n}\ket{0}^{\otimes n}.
\end{equation}
Al medir los $n$ qubits se lee directamente $a$. \textbf{Una sola consulta}, es
decir, complejidad $O(1)$ frente a $O(n)$ clásica.

\paragraph{Sin entrelazamiento.} El oráculo se arma como un producto de $n$
puertas de un solo qubit: $U_i=I$ si $a_i=0$ y $U_i=Z$ si $a_i=1$. Todas las
puertas del circuito son de un qubit y el estado inicial no está entrelazado,
así que el estado nunca lo está. La aceleración viene de la
\textbf{superposición}, no del entrelazamiento.

\subsection{Código (completamente comentado)}
\lstinputlisting[style=py]{../../fuentes/16_bernstein_vazirani.py}

\subsection{Explicación celda por celda}
\begin{description}
  \item[Celda 2] Define el número oculto $a=1011$ y la caja negra $f(x)=a\cdot x$.
  \item[Celda 3] Solución clásica: cuatro consultas con $x=1000,0100,0010,0001$.
  \item[Celda 4] El oráculo: una $Z$ en cada qubit donde $a_i=1$.
  \item[Celda 5] El circuito $H^{\otimes n}$, $U_f$, $H^{\otimes n}$ y medición.
  \item[Celda 6] Lo ejecuta (1000 mediciones) e invierte la cadena, porque
    Qiskit escribe el qubit~0 a la derecha.
  \item[Celda 7] Calcula la entropía de cada qubit antes de medir para
    comprobar que no hay entrelazamiento.
\end{description}

\subsection{Salida obtenida}
\lstinputlisting[style=salida]{../../fuentes/salidas/16_bernstein_vazirani.txt}
(El dibujo del circuito no se reproduce aquí.)

\subsection{Interpretación}
Las cuatro consultas clásicas dan $1,0,1,1$, es decir, $a=1011$. El circuito
cuántico, con una sola consulta, mide siempre \texttt{1101} en 1000 mediciones;
al invertir la cadena se lee \texttt{1011}, el número oculto. La entropía de
cada qubit es $0$: el estado nunca está entrelazado.

\subsection{Ejercicios}
\begin{enumerate}
  \item Prueba con otro número oculto, por ejemplo $a=0110$ y $a=11111$.
  \item ¿Qué pasa si $a=0000$? Predícelo y compruébalo.
  \item Cuenta cuántas consultas necesita el algoritmo clásico para
    $n=20$ y cuántas el cuántico.
\end{enumerate}

\subsection{Errores comunes}
\begin{itemize}
  \item Olvidar que Qiskit imprime el qubit~0 a la derecha: la cadena
    medida está al revés respecto de $a_1a_2\ldots a_n$.
  \item Confundir el número de \emph{consultas} con el tiempo total: el
    algoritmo cuenta cuántas veces se usa la caja negra.
\end{itemize}

\ifdefined\enmain\else
  \end{document}
\fi

